\chapter{Martingale Maximal Inequalities and Convergence in Mean Power}
\label{sec:p15-martingale-maximal-inequalities}

A maximal inequality bounds an entire finite history by information at
its terminal time. Doob's inequalities turn this control into convergence
in \(L^p\), while their endpoint counterexamples show why a first-moment
bound is insufficient. The source presentation is Durrett's complete
Section~4.4, including its exercises and branching-process application
\cite[pp.~235--239]{DurrettPTEAuthor2019}.
The branching prerequisites from Section~4.3.4 are supplied below.

\section{Stopping and finite-history maxima}
\label{sec:p15-maximal-stopping}

\begin{hbdefinition}{def:p15-martingale-maxima}{Positive and absolute maxima}
All processes in this chapter are real and adapted to a specified
filtration \((\mathcal F_n)_{n\geq0}\) on \((\Omega,\mathcal F,P)\).
For a process \(X\), write
\[
  \overline X_n=\max_{0\leq j\leq n}X_j^+,\qquad
  X_n^*=\max_{0\leq j\leq n}|X_j|,\qquad
  X_\infty^*=\sup_{j\geq0}|X_j|.
\]
The finite maxima and the countable supremum are measurable.
For \(1\leq p<\infty\), use
\(\|Y\|_p=(E|Y|^p)^{1/p}\), permitting an infinite value.
The symbol \(\|Y\|_\infty\) denotes the smallest extended bound \(c\)
with \(|Y|\leq c\) almost surely.
For an almost-surely finite stopping time \(T\), define \(X_T=0\)
on the null event \(T=\infty\); on \(T=j<\infty\) use \(X_j\).
This convention supplies a real measurable version, not a limit at infinity.
\index{Doob maximal inequality}\index{martingale!maximum}
\end{hbdefinition}

\begin{hbtheorem}{thm:p15-bounded-stopping-routes}{proved=annex:p15.bounded-stopping-routes}{The two source stopping proofs and conditional sampling}
If \(X\) is a submartingale, \(T\) is a stopping time, and
\(P(T\leq k)=1\) for an integer
\(k\geq0\), then
\[
                         EX_0\leq EX_T\leq EX_k.
\]
The annex supplies both the complementary predictable-transform proof
and the event-by-event proof requested by the source.

If stopping times satisfy \(S\leq T\) pointwise and \(A\in\mathcal F_S\),
then \(L=S\) on \(A\), \(L=T\) on \(A^c\), is a stopping time.
If also \(P(T\leq k)=1\), then \(X_S,X_T\) are integrable and
\[
                      X_S\leq E(X_T\mid\mathcal F_S),\qquad
                      EX_S\leq EX_T.
\]
For supermartingales the inequalities reverse; for martingales they
are equalities \cite[thm.~4.4.1, exs.~4.4.1--4]{DurrettPTEAuthor2019}.

The almost-sure bound on \(T\) is a hypothesis of the first inequality,
not a convenience. For the walk of \cref{thm:p15-fair-walk-maximal-tail}
started at \(S_0=1\), with \(T_0=\inf\{n\geq0:S_n=0\}\), the walk is a
martingale, hence a submartingale, and \(T_0\) is an almost-surely finite
stopping time that is not almost surely bounded; yet
\[
                    ES_0=1>0=ES_{T_0},
\]
so \(EX_0\leq EX_T\) fails for this unbounded \(T\)
\cite[remark after thm.~4.4.1, p.~235]{DurrettPTEAuthor2019}.
The same remark points ahead to its Section~5.7 for conditions under
which \(EX_0\leq EX_N\) holds for unbounded \(N\). The source's
Section~4.8 opens with this very question, citing the same walk, and
answers it with uniform-integrability conditions
\cite[sec.~4.8, thms.~4.8.2--4.8.3, pp.~255--256]{DurrettPTEAuthor2019};
those conditions are proved in
\cref{thm:p15-os-ui-stopped,thm:p15-os-residual-ui}.
\end{hbtheorem}

\begin{hbtheorem}{thm:p15-doob-weak-maximal}{proved=annex:p15.doob-weak-maximal}{Doob's probability inequality}
For a submartingale \(X\), a finite \(n\geq0\) and \(\lambda>0\), put
\(A=\{\overline X_n\geq\lambda\}\). Then
\[
                         \lambda P(A)\leq E(X_n;A)\leq EX_n^+.
\]
The same inequality holds with \(A=\{\overline X_n>\lambda\}\).
For a martingale \(Y\), applying it to \(|Y|\) gives
\[
                  \lambda P(Y_n^*\geq\lambda)
                         \leq E(|Y_n|;Y_n^*\geq\lambda)
                         \leq E|Y_n|.
\]
Time zero is included in each maximum
\cite[thm.~4.4.2]{DurrettPTEAuthor2019}.
\end{hbtheorem}

\begin{hbtheorem}{thm:p15-kolmogorov-martingale-route}{proved=annex:p15.kolmogorov-martingale-route}{The separate martingale proof of Kolmogorov's inequality}
Let \(\xi_i\) be independent real variables with
\(E\xi_i=0\) and \(E\xi_i^2=\sigma_i^2<\infty\).
Put \(S_0=0\), \(S_n=\sum_{i=1}^n\xi_i\) and
\(v_n=\sum_{i=1}^n\sigma_i^2\). For \(x>0\),
\[
               P\left(\max_{0\leq m\leq n}|S_m|\geq x\right)
                                      \leq\frac{v_n}{x^2}.
\]
This proof uses the square submartingale and Doob's inequality,
retaining the source's distinct route to the earlier independent-sum
bound \cite[example~4.4.3, p.~235]{DurrettPTEAuthor2019}.
\end{hbtheorem}

\begin{hbtheorem}{thm:p15-doob-lp-maximal}{proved=annex:p15.doob-lp-maximal}{The maximal \(L^p\) inequality}
For a submartingale \(X\) and \(1<p<\infty\),
\[
                 E(\overline X_n)^p
                    \leq\left(\frac p{p-1}\right)^pE(X_n^+)^p.
\]
For a martingale \(Y\),
\[
                 E(Y_n^*)^p
                    \leq\left(\frac p{p-1}\right)^pE|Y_n|^p.
\]
The inequalities are valid as extended nonnegative inequalities.
The proof first truncates the maximum, so it never divides by an
unproved finite moment or by zero
\cite[thm.~4.4.4]{DurrettPTEAuthor2019}.
\end{hbtheorem}

\begin{hbtheorem}{thm:p15-fair-walk-maximal-tail}{proved=annex:p15.fair-walk-maximal-tail}{The exact fair-walk endpoint counterexample}
Let \(S_0=1\), with independent increments \(1,-1\) of probability
\(1/2\), and let \(T_0=\inf\{n\geq0:S_n=0\}\).
The process \(X_n=S_{n\wedge T_0}\) is a nonnegative martingale,
\(EX_n=1\), and \(T_0<\infty\) almost surely. For every integer \(m\geq1\),
\[
                 P\left(\sup_{n\geq0}X_n\geq m\right)=\frac1m.
\]
Consequently \(E\sup_nX_n=\infty\), while
\(E\max_{j\leq n}X_j\uparrow\infty\) and \(EX_n=1\).
There is no finite universal constant bounding the expected maximum
by the terminal first absolute moment. The hitting calculation is
proved using bounded stopping and bounded convergence
\cite[example~4.4.5, pp.~236--237]{DurrettPTEAuthor2019}.
\end{hbtheorem}

\begin{hbtheorem}{thm:p15-martingale-lp-convergence}{proved=annex:p15.martingale-lp-convergence}{Convergence under a bounded \(p\)th moment}
Let \(1<p<\infty\). If \(X\) is a martingale and
\(C=\sup_nE|X_n|^p<\infty\), then it has a real integrable limit
\(X_\infty\), almost surely and in \(L^p\), with
\[
      E(X_\infty^*)^p\leq\left(\frac p{p-1}\right)^pC,\qquad
      E|X_\infty|^p\leq C,\qquad
      X_n=E(X_\infty\mid\mathcal F_n).
\]
The notation \(X_\infty^*\) still denotes the supremum of the original
sequence. A nonnegative submartingale with bounded \(p\)th moments
also converges almost surely and in \(L^p\).

The source states the martingale case, convergence almost surely and
in \(L^p\) for \(p>1\), and obtains the integrability of the supremum
inside its proof \cite[thm.~4.4.6, p.~237]{DurrettPTEAuthor2019}.
The displayed bound with the constant \(C\), the bound
\(E|X_\infty|^p\leq C\), the closure identity and the
nonnegative-submartingale extension are supplied here; the source
states none of them in its Section~4.4.
Neither hypothesis can be dropped, and the examples showing this are
also supplied here rather than taken from the source.
\Cref{thm:p15-maximal-convergence-boundaries}~(1) gives, for each
\(p>1\), a submartingale taking negative values with
\(E|X_n|^p=1\) for every \(n\) that converges almost surely but not in
\(L^p\). \Cref{thm:p15-maximal-convergence-boundaries}~(2) gives a
martingale bounded by \(1\), hence with bounded \(p\)th moments for
every finite \(p\) and bounded uniform norms, whose limit is
attained almost surely and in every finite \(L^p\) but not in
\(L^\infty\).
\end{hbtheorem}

\section{Orthogonality and mean-square identities}
\label{sec:p15-martingale-l2-identities}

\begin{hbtheorem}{thm:p15-martingale-orthogonality}{proved=annex:p15.martingale-orthogonality}{Orthogonality to earlier information}
Let \(X\) be a martingale with \(EX_n^2<\infty\) for every \(n\).
If \(m\leq n\) and \(Y\) is \(\mathcal F_m\)-measurable with
\(EY^2<\infty\), then
\[
                            E((X_n-X_m)Y)=0.
\]
In particular distinct increments are orthogonal in \(L^2\), and
\[
       E(X_n-X_0)^2=\sum_{j=1}^nE(\Delta X_j)^2,\qquad
       EX_n^2=EX_0^2+\sum_{j=1}^nE(\Delta X_j)^2.
\]
Orthogonality does not assert independence
\cite[thm.~4.4.7]{DurrettPTEAuthor2019}.
\end{hbtheorem}

\begin{hbtheorem}{thm:p15-martingale-conditional-variance}{proved=annex:p15.martingale-conditional-variance}{Conditional variance and nested projections}
For the same martingale and \(m\leq n\),
\[
             E((X_n-X_m)^2\mid\mathcal F_m)
                         =E(X_n^2\mid\mathcal F_m)-X_m^2.
\]
If \(\mathcal H\subseteq\mathcal G\) and \(Y\in L^2\), then
\[
 E\bigl(E(Y\mid\mathcal G)-E(Y\mid\mathcal H)\bigr)^2
   =E\bigl(E(Y\mid\mathcal G)\bigr)^2
                      -E\bigl(E(Y\mid\mathcal H)\bigr)^2.
\]
Both formulas include all square-integrability and nested-information
conditions \cite[thm.~4.4.8, ex.~4.4.5]{DurrettPTEAuthor2019}.
The first formula is the conditional counterpart of the variance identity
\(E(X-EX)^2=EX^2-(EX)^2\): the constant \(EX\) is replaced by the
\(\mathcal F_m\)-measurable best prediction \(X_m=E(X_n\mid\mathcal F_m)\),
and expectation by conditional expectation given \(\mathcal F_m\).
The source remarks that it is proved in exactly the same way
\cite[remark after thm.~4.4.8, p.~237]{DurrettPTEAuthor2019}, and the annex
writes both computations side by side.
\end{hbtheorem}

\begin{hbtheorem}{thm:p15-martingale-covariance-isometry}{proved=annex:p15.martingale-covariance-isometry}{Covariances and predictable weighted sums}
For square-integrable martingales \(X,Y\) on the same filtration,
\[
        E(X_nY_n)-E(X_0Y_0)
                        =\sum_{j=1}^nE(\Delta X_j\Delta Y_j).
\]
If \(H_j,K_j\) are predictable and bounded by deterministic constants
at each finite time, then
\[
 E\left[\left(\sum_{j=1}^nH_j\Delta X_j\right)
         \left(\sum_{j=1}^nK_j\Delta Y_j\right)\right]
                    =\sum_{j=1}^nE(H_jK_j\Delta X_j\Delta Y_j).
\]
In particular
\[
                 E\left(\sum_{j=1}^nH_j\Delta X_j\right)^2
                                =\sum_{j=1}^nE(H_j^2(\Delta X_j)^2).
\]
These formulas require no independence of the increments
\cite[ex.~4.4.9]{DurrettPTEAuthor2019}.
\end{hbtheorem}

\begin{hbtheorem}{thm:p15-square-summable-martingale}{proved=annex:p15.square-summable-martingale}{Square-summable increments give a complete limit}
For a martingale with
\[
                       EX_0^2+\sum_{j=1}^\infty E(\Delta X_j)^2<\infty,
\]
there is a limit \(X_\infty\), almost surely and in \(L^2\), and
\[
                  E(X_\infty-X_n)^2
                                =\sum_{j>n}E(\Delta X_j)^2.
\]
Thus the finite partial-sum process of square-integrable martingale
differences converges when their second moments are summable
\cite[ex.~4.4.10]{DurrettPTEAuthor2019}.
\end{hbtheorem}

\begin{hbtheorem}{thm:p15-weighted-martingale-law}{proved=annex:p15.weighted-martingale-law}{The source's weighted martingale strong law}
Let \(X\) be a martingale, with \(E|X_n|<\infty\) for every \(n\) and no
moment of \(X_0\) beyond its integrability assumed, and write
\(\xi_m=X_m-X_{m-1}=\Delta X_m\) for \(m\geq1\). Let \((b_m)_{m\geq1}\) be
nondecreasing with \(b_m\to\infty\) and every \(b_m\neq0\), so that each
term of the series below is defined. If
\[
                           \sum_{m=1}^\infty
                                  \frac{E\xi_m^2}{b_m^2}<\infty,
\]
then \(X_n/b_n\to0\) almost surely.
In particular a uniform bound \(E\xi_n^2\leq K<\infty\), together
with \(\sum_mb_m^{-2}<\infty\), suffices.
The source states this as Exercise~4.4.11, which continues only the
notation \(\xi_m\) of Exercise~4.4.10; the hypothesis \(EX_0^2<\infty\)
of Exercise~4.4.10 is not part of it
\cite[exs.~4.4.10--4.4.11, p.~239]{DurrettPTEAuthor2019}. The statement
therefore covers martingales whose initial value is integrable but not
square integrable, for which no \(X_n\) is square integrable. For
example, let \(X_0\geq1\) have density \(2x^{-3}\) on \([1,\infty)\), so
that \(EX_0=2\) and \(EX_0^2=\infty\), let \(\xi_1,\xi_2,\ldots\) be
independent fair signs independent of \(X_0\), and put
\(X_n=X_0+\xi_1+\cdots+\xi_n\) and \(b_m=m\). The series is
\(\sum_mm^{-2}<\infty\), so \(X_n/n\to0\) almost surely, although
\(EX_n^2=\infty\) for every \(n\).
The proof constructs the weighted martingale and then applies the
proved Kronecker lemma.
\end{hbtheorem}

\section{Further maximal bounds and their endpoints}
\label{sec:p15-maximal-refinements}

\begin{hbtheorem}{thm:p15-bounded-increment-small-ball}{proved=annex:p15.bounded-increment-small-ball}{A small-ball bound with bounded overshoot}
For the independent centered sums of
\cref{thm:p15-kolmogorov-martingale-route}, suppose
\(|\xi_j|\leq K\) almost surely. For \(x\geq0\) and \(v_n>0\),
\[
                 P\left(\max_{1\leq j\leq n}|S_j|\leq x\right)
                                      \leq\frac{(x+K)^2}{v_n}.
\]
The exit time uses the strict condition \(|S_j|>x\), preserving
the weak inequality inside the probability
\cite[ex.~4.4.6]{DurrettPTEAuthor2019}.
The proof uses that \(S_j^2-v_j\) is a martingale when the increments
are independent and centered with possibly different variances.
The exercise takes this from its earlier Exercise~4.2.2; in the cited
edition that exercise asks for a different example, and the nearest
statement, Example~4.2.2, treats identically distributed increments
\cite[example~4.2.2, p.~218; ex.~4.2.2, p.~223]{DurrettPTEAuthor2019}.
The annex therefore proves the conditional square identity
\(E(S_{j+1}^2\mid\mathcal F_j)=S_j^2+\sigma_{j+1}^2\) directly.
\end{hbtheorem}

\begin{hbtheorem}{thm:p15-cantelli-maximal}{proved=annex:p15.cantelli-maximal}{An optimized one-sided quadratic bound}
If \(X_0=0\) and \(X\) is a square-integrable martingale, then for
\(\lambda>0\), with \(v=EX_n^2\),
\[
                   P\left(\max_{0\leq j\leq n}X_j\geq\lambda\right)
                                      \leq\frac{v}{v+\lambda^2}.
\]
For any prescribed \(v>0,\lambda>0\), a one-step two-point martingale
attains this bound. The case \(v=0\) gives probability zero
\cite[ex.~4.4.7]{DurrettPTEAuthor2019}.

The source introduces this exercise with the sentence
``The next result gives an extension of Theorem~4.4.2 to \(p=1\)''
\cite[ex.~4.4.7, p.~239]{DurrettPTEAuthor2019}.
What the result verifiably adds to Doob's probability inequality
(\cref{thm:p15-doob-weak-maximal}, the source's Theorem~4.4.2) is the
following, each part proved in the annex. It is obtained from that
inequality itself, applied to the submartingales \((X_j+c)^2\).
Applied directly to the submartingale \((X_j^+)^2\), Doob's inequality
gives only \(P(\max_{j\leq n}X_j\geq\lambda)\leq v/\lambda^2\); for
\(v>0\) the bound above is strictly smaller,
\(v/(v+\lambda^2)<\min(1,v/\lambda^2)\), it is never trivial, and the
two-point law shows that it cannot be lowered. The hypotheses and the
bound are second-moment quantities, and the handbook does not identify
the exponent \(1\) in the source's phrase with any exponent in the
statement; the phrase is reported here as the source's framing, not as
part of the result.
\end{hbtheorem}

\begin{hbtheorem}{thm:p15-doob-llogl}{proved=annex:p15.doob-llogl}{The logarithmic substitute at the first-moment endpoint}
Set \(\log^+0=0\) and \(t\log^+t=0\) at \(t=0\), and
\(\log^+t=\max(0,\log t)\) for \(t>0\).
For a submartingale \(X\),
\[
      E\overline X_n
        \leq\frac1{1-e^{-1}}
                  \left[1+E\bigl(X_n^+\log^+X_n^+\bigr)\right].
\]
More generally, for every \(a>0\),
\[
      E\overline X_n
        \leq\frac1{1-e^{-1}}
                  \left[a+E\bigl(X_n^+\log^+(X_n^+/a)\bigr)\right].
\]
The truncated-maximum proof uses \(\log^+\), so no logarithm of a
zero maximum is introduced \cite[ex.~4.4.8]{DurrettPTEAuthor2019}.
\end{hbtheorem}

\section{The complete branching-process application}
\label{sec:p15-branching-second-moments}

\begin{hbdefinition}{def:p15-branching-genealogy}{Offspring variables and independent descendant families}
Let \(\mathcal T=\bigcup_{r\geq0}\mathbb N_{>0}^r\) be the countable
set of finite sequences of positive integers, with empty sequence
\(\varnothing\). On its countable product space assign independent
offspring variables \((\xi_v)_{v\in\mathcal T}\) with a common
nonnegative integer law \(p_k=P(\xi_v=k)\).
The realized root set is \(V_0=\{\varnothing\}\), and
\[
              V_{n+1}=\{vi:v\in V_n,\ 1\leq i\leq\xi_v\},
              \qquad Z_n=\operatorname{card}(V_n).
\]
Write \(|v|_{\mathrm{len}}\) for sequence length and use the filtration
\(\mathcal H_n=\sigma(\xi_v:|v|_{\mathrm{len}}<n)\).
Every generation is finite. Put
\[
              \mu=\sum_{k\geq0}kp_k<\infty,\qquad
              f(s)=\sum_{k\geq0}p_ks^k\quad(0\leq s\leq1),
\]
where \(s^0=1\), including at \(s=0\).
When \(\mu>0\), define \(W_n=Z_n/\mu^n\).
For an initial population \(z\in\mathbb N_0\), use \(z\) independent
copies of this tree and sum their generation counts.
The process \((Z_n)_{n\geq0}\) is a \emph{Galton--Watson process}, and
\((p_k)\) is its \emph{offspring distribution}. The source gives both
names to the process of its recursion
\(Z_{n+1}=\xi_1^{\,n+1}+\cdots+\xi_{Z_n}^{\,n+1}\) when \(Z_n>0\) and
\(Z_{n+1}=0\) when \(Z_n=0\), with \(Z_0=1\) and i.i.d.\ nonnegative
integer variables \(\xi_i^{\,n}\) \cite[p.~230, eq.~(4.3.4)]{DurrettPTEAuthor2019};
by \cref{thm:p15-branching-family-construction} the tree counts above
have the same law. In either form \(Z_n\) is the number of individuals
in generation \(n\), and each individual of a generation has,
independently of the others, a number of children with law \((p_k)\).
\index{Galton--Watson process}\index{branching process!second moment}
\end{hbdefinition}

\begin{hbtheorem}{thm:p15-branching-family-construction}{proved=annex:p15.branching-family-construction}{Generation laws, martingales and independent subtrees}
The tree construction has the same generation-count law as the source's
recursion \(Z_{n+1}=\sum_{i=1}^{Z_n}\xi_i^{\,n+1}\), where the
array variables are i.i.d. with the same offspring law. It satisfies
\[
                 E(Z_{n+1}\mid\mathcal H_n)=\mu Z_n,\qquad
                 EZ_n=\mu^n\quad(n\geq1).
\]
Here \(EZ_0=1\) also when \(\mu=0\).
For \(\mu>0\), \(W_n\) is a nonnegative martingale with mean one
and a finite almost-sure limit \(W\).
The normalized row process is likewise a martingale for the source
filtration \(\mathcal F_n^{\mathrm{row}}
=\sigma(\xi_i^m:i\geq1,\ 1\leq m\leq n)\).
The descendants of different children are independent copies of the
whole process, independent of the root offspring count.
These facts supply the source's Lemma~4.3.9 and the family decomposition
used in its later proofs \cite[pp.~230--232]{DurrettPTEAuthor2019}.
\end{hbtheorem}

\begin{hbtheorem}{thm:p15-branching-extinction-root}{proved=annex:p15.branching-extinction-root}{Extinction and the exact generating-function root}
Let \(D=\bigcup_n\{Z_n=0\}\) and \(\rho=P(D)\).
With \(\theta_0=0\), \(\theta_n=P(Z_n=0)\), one has
\[
                 \theta_{n+1}=f(\theta_n),\qquad
                 \theta_n\uparrow\rho,
\]
and \(\rho\) is the smallest solution of \(f(s)=s\) in \([0,1]\).
\begin{enumerate}
\item If \(0\leq\mu<1\), extinction is almost sure.
\item If \(\mu=1\) and \(p_1<1\), extinction is almost sure.
      If \(p_1=1\), then \(Z_n=1\) for all \(n\).
\item If \(\mu>1\), there is exactly one root in \([0,1)\),
      and it is \(\rho<1\).
\end{enumerate}
The derivative, convexity and root arguments are proved, including
the endpoint \(s=1\)
\cite[thms.~4.3.10--12]{DurrettPTEAuthor2019}.

In the extinction cases the normalized process has limit zero.
If \(0<\mu<1\), or \(\mu=1\) and \(p_1<1\), then almost surely
\(Z_n=0\) for all sufficiently large \(n\), and therefore
\[
                 W_n=\frac{Z_n}{\mu^n}\longrightarrow0
                 \quad\text{almost surely};
\]
the limit is zero because the population dies out, not because of
the normalization. For \(0<\mu<1\) the probability of survival to
generation \(n\) is at most \(\mu^n\), so it decays exponentially fast
\cite[thm.~4.3.10, p.~231; p.~232]{DurrettPTEAuthor2019}.
The excluded critical case \(p_1=1\) has \(W_n=Z_n=1\) for all \(n\),
so its limit is \(1\); when \(\mu=0\) the normalization is undefined and
\(Z_1=0\) already. The source calls the subcritical conclusion intuitive:
if each individual has on average fewer than one child, the species
dies out \cite[p.~231]{DurrettPTEAuthor2019}. The proof makes that
reading exact through the mean bound \(P(Z_n>0)\leq EZ_n=\mu^n\).
\end{hbtheorem}

\begin{hbtheorem}{thm:p15-branching-second-moment}{proved=annex:p15.branching-second-moment}{The full supercritical \(L^2\) calculation}
Assume \(\mu>1\) and
\(\sigma^2=\operatorname{Var}(\xi_v)<\infty\), with \(Z_0=1\).
Then
\[
 \begin{split}
 E((W_n-W_{n-1})^2\mid\mathcal H_{n-1})
                         &=\sigma^2Z_{n-1}\mu^{-2n},\\
 EW_n^2&=1+\sigma^2\sum_{j=1}^n\mu^{-j-1}
        =1+\frac{\sigma^2}{\mu(\mu-1)}(1-\mu^{-n}).
 \end{split}
\]
The limit exists almost surely and in \(L^2\), with
\[
                    EW=1,\qquad
                    EW^2=1+\frac{\sigma^2}{\mu(\mu-1)}.
\]
In particular \(P(W>0)>0\)
\cite[example~4.4.9, p.~238]{DurrettPTEAuthor2019}.
\end{hbtheorem}

\begin{hbtheorem}{thm:p15-branching-positive-limit}{proved=annex:p15.branching-positive-limit}{Positivity exactly on survival}
For any finite \(\mu>1\), if the martingale limit \(W\) is not
almost surely zero, then
\[
                 P(W=0)=\rho,\qquad
                 \{W>0\}=\{Z_n>0\text{ for every }n\}
                                      \quad\text{almost surely}.
\]
This applies in particular under the finite-variance hypotheses above.
For initial population \(z\in\mathbb N_0\), extinction probability is
\(\rho^z\), with \(\rho^0=1\). The process \(\rho^{Z_n}\) is a
bounded martingale; its convergence also gives this extinction formula
\cite[exs.~4.3.11--12]{DurrettPTEAuthor2019}.
\end{hbtheorem}

\begin{hbtheorem}{thm:p15-branching-calculations}{proved=annex:p15.branching-calculations}{The source's binomial example and limiting cases}
For offspring law \(\operatorname{Binomial}(3,1/2)\),
\[
   f(s)=\left(\frac{1+s}{2}\right)^3,\quad
   \mu=\frac32,\quad \sigma^2=\frac34,\quad
   \rho=\sqrt5-2,\quad P(W>0)=3-\sqrt5,\quad EW^2=2.
\]
The iterates \(0,f(0),f(f(0)),\ldots\) increase to this root.
Their first positive values are \(\theta_1=1/8\) and
\(\theta_2=729/4096\).
The binomial family-name law is the source's
\cite[ex.~4.3.13 and figs.~4.3--4.4]{DurrettPTEAuthor2019}.
The source introduces it through the question that, on its account,
led Galton and Watson to the process named after them: whether a family
name survives. Each family is assumed to have exactly three children,
each child's sex is fixed by an independent fair coin flip, and, as in
the nineteenth century, only sons pass on the name. Following the male
line then gives a branching process with
\(p_0=1/8\), \(p_1=3/8\), \(p_2=3/8\), \(p_3=1/8\), which is the law
\(\operatorname{Binomial}(3,1/2)\), and the exercise asks for the
probability \(\rho\) that the name dies out when \(Z_0=1\)
\cite[ex.~4.3.13, p.~234]{DurrettPTEAuthor2019}. The value
\(\rho=\sqrt5-2\) above answers it: under these assumptions a single
founder's name dies out with probability about \(0.236\).
For deterministic offspring \(b\geq2\), \(W_n=1\) and \(\rho=0\).
For \(\operatorname{Binomial}(2,1/2)\), \(Z_n\to0\) almost surely
but \(EZ_n=1\), so convergence is not in \(L^1\).
These two limiting offspring laws are supplied here; they do not occur
in the source. They are instances of its supercritical root theorem
(\(\mu=b>1\), with \(f(s)=s^b\) and root \(0\)) and of its critical
extinction theorem (\(\mu=1\), \(p_1=1/2<1\)), whose mean identity
\(EZ_n=\mu^n\) is part of the proof of its subcritical theorem
\cite[thms.~4.3.10--4.3.12, pp.~231--232]{DurrettPTEAuthor2019}.
\end{hbtheorem}

\section{Why the convergence hypotheses matter}
\label{sec:p15-maximal-convergence-boundaries}

\begin{hbtheorem}{thm:p15-maximal-convergence-boundaries}{proved=annex:p15.maximal-convergence-boundaries}{Signed submartingales, uniform norm and independence}
\begin{enumerate}
\item Fix \(p>1\). On \((0,1]\) with uniform probability put
      \(A_n=(0,1/(n+1)]\), \(\mathcal F_n=\sigma(A_0,\ldots,A_n)\), and
      \(X_n=-(n+1)^{1/p}\mathbf1_{A_n}\).
      This is a submartingale with \(E|X_n|^p=1\) and \(X_n\to0\)
      almost surely, but it does not converge in \(L^p\).
      Moreover \(E(X_n^*)^p=\sum_{j=1}^{n+1}1/j\).
\item On the dyadic probability space of
      \cref{def:p15-dyadic-information}, let
      \(X_n=E(\mathbf1_{[0,1/3)}\mid\mathcal D_n)\).
      This bounded martingale converges almost surely and in every
      finite \(L^p\), but
      \(\|X_n-\mathbf1_{[0,1/3)}\|_\infty=2/3\) for every \(n\).
\item For independent fair signs \(\varepsilon_1,\varepsilon_2\),
      the increments
      \(\xi_1=\varepsilon_1\), \(\xi_2=(1+\varepsilon_1)\varepsilon_2\)
      form martingale differences in their input filtration.
      They are orthogonal but dependent, since
      \(P(\xi_2=0\mid\varepsilon_1=-1)=1\) and \(P(\xi_2=0)=1/2\).
\end{enumerate}
\end{hbtheorem}

The predictable square-variation construction and its convergence,
normalization and stopping estimates are supplied in
\cref{thm:p15-predictable-variation-identities,thm:p15-finite-variation-convergence,thm:p15-variance-clock-normalization,thm:p15-square-root-variation-bound,thm:p15-square-root-stopping}.
The finite-variance branching result does not replace the
separate Kesten--Stigum logarithmic-moment criterion referred to by the
source in Theorem~4.3.13.
